\textbf{COLUNI 2020 -- Questão 32}\par\medskip Leia o texto. Fenômeno dos rios voadores

\par

“Os rios voadores são ‘cursos de água atmosféricos’, formados por massas de ar carregadas de vapor de água, muitas vezes acompanhados por nuvens, e são propelidos pelos ventos. Essas correntes de ar invisíveis passam por cima de nossas cabeças carregando umidade da Bacia Amazônica para o Centro- Oeste, Sudeste e Sul do Brasil.” Disponível em: https://rios.voadores.com.br. Acesso em: 18 set. 2019.

\par

Com relação aos rios voadores, analise as ocorrências apresentadas abaixo e, a seguir, numere-as de acordo com a ordem cronológica em que acontecem.

\par

(   ) A água que se acumula no solo e nas árvores evapora, formando as nuvens. (   ) A água evapora do Oceano Atlântico e forma nuvens que são levadas para a floresta pelos ventos. (   ) A barreira geográfica dos Andes redireciona os ventos para o centro do continente. (   ) As nuvens carregadas fazem chover sobre a floresta. (   ) As novas nuvens vão para o Centro-Oeste, Sudeste e Sul, fazendo chover nessas regiões e, posteriormente, seguem viagem pela América do Sul.

\par

Assinale a sequência CORRETA:\begin{enumerate}
\item 1, 2, 3, 4, 5
\item 2, 4, 1, 3, 5
\item 3, 1, 4, 2, 5
\item 4 ,1, 3, 2, 5\end{enumerate}
